\subsection{Capture Timeout Risk}
\label{sec:capture-timeout-risk}
\label{sec:guard-limit}

Enabling lightweight multi-frame composition in the draft sequence increases the risk of a Capture Timeout, especially under parallel capture.
Since each capture's timeout clock starts at request acceptance, waiting in the draft queue consumes part of its timeout budget before its draft sequence starts.
Figure~\ref{fig:overlap} illustrates how this waiting time grows across successive captures even when their draft sequences take equally long.
The additional processing required for lightweight multi-frame composition can lengthen these waits and cause later captures to miss their deadlines.
Because a single Capture Timeout violation blocks release, Samsung has not deployed this feature on any device model.

\begin{figure}[t]
  \centering
  \resizebox{\columnwidth}{!}{% Notes: docs/exhibits.md#fig_parallel_capture_overlap
\begin{tikzpicture}[font=\small,
  collectstage/.style={draw=black!55, fill=gray!14},
  draftstage/.style={draw=blue!45!black, fill=blue!14},
  queued/.style={pattern=north east lines, pattern color=red!45, draw=red!35},
  overrun/.style={draw=red!80!black, fill=red!65},
  ddl/.style={red, dashed, thick},
  sig/.style={->, black!75}]

  \node[anchor=east] at (-0.2, 3.925) {capture $i$};
  \draw[very thick] (0, 3.55) -- (0, 4.30);
  \draw[collectstage] (0.05, 3.70) rectangle (0.90, 4.15);
  \node at (0.475, 3.925) {\scriptsize collect};
  \draw[draftstage] (0.90, 3.70) rectangle (3.50, 4.15);
  \node at (2.20, 3.925) {\scriptsize Draft Sequence};

  \draw[densely dotted, gray] (3.50, 3.925) -- (7.30, 3.925);
  \node[fill=white, inner sep=1.5pt] at (5.40, 3.925)
    {\footnotesize\textcolor{green!45!black}{\ding{51} in time}};
  \draw[ddl] (7.30, 3.55) -- (7.30, 4.30);
  \node[red, anchor=east] at (9.00, 3.925) {deadline $i$};

  \fill[black!60] (0.35, 3.68) circle (1.3pt);
  \draw[sig] (0.35, 3.66) -- (0.50, 3.17);
  \node[anchor=west, black!75] at (0.42, 3.42)
    {\scriptsize\texttt{captureAvailable}};

  \node[anchor=east] at (0.3, 2.725) {capture $i{+}1$};
  \draw[very thick] (0.50, 2.35) -- (0.50, 3.10);
  \draw[collectstage] (0.55, 2.50) rectangle (1.40, 2.95);
  \node at (0.975, 2.725) {\scriptsize collect};
  \draw[queued] (1.40, 2.50) rectangle (3.50, 2.95);
  \node[fill=white, inner sep=1.5pt] at (2.45, 2.725)
    {\footnotesize\textcolor{red!60!black}{waiting}};
  \draw[draftstage] (3.50, 2.50) rectangle (6.10, 2.95);
  \node at (4.80, 2.725) {\scriptsize Draft Sequence};

  \draw[densely dotted, gray] (6.10, 2.725) -- (7.80, 2.725);
  \node[fill=white, inner sep=1.5pt] at (6.95, 2.725)
    {\footnotesize\textcolor{green!45!black}{\ding{51} in time}};
  \draw[ddl] (7.80, 2.35) -- (7.80, 3.10);
  \node[red, anchor=east] at (10.00, 2.725) {deadline $i+1$};

  \fill[black!60] (0.85, 2.48) circle (1.3pt);
  \draw[sig] (0.85, 2.46) -- (1.00, 1.97);

  \node[anchor=east] at (0.8, 1.525) {capture $i{+}2$};
  \draw[very thick] (1.00, 1.15) -- (1.00, 1.90);
  \draw[collectstage] (1.05, 1.30) rectangle (1.90, 1.75);
  \node at (1.475, 1.525) {\scriptsize collect};
  \draw[queued] (1.90, 1.30) rectangle (6.10, 1.75);
  \node[fill=white, inner sep=1.5pt] at (4.00, 1.525)
    {\footnotesize\textcolor{red!60!black}{waiting}};
  \draw[draftstage] (6.10, 1.30) rectangle (8.70, 1.75);

  \draw[overrun] (8.30, 1.30) rectangle (8.70, 1.75);
  \node at (7.20, 1.525) {\scriptsize Draft Sequence};
  \draw[ddl] (8.30, 1.15) -- (8.30, 1.90);
  \node[red!70!black, anchor=west, align=left] at (8.80, 1.525)
    {\footnotesize\bfseries Capture\\[-2pt]\footnotesize\bfseries Timeout};

  \draw[dashed, gray, ->] (3.50, 3.67) -- (3.50, 2.98);
  \draw[dashed, gray, ->] (6.10, 2.47) -- (6.10, 1.78);

  \draw[decorate, decoration={brace, mirror, amplitude=4pt}]
    (1.00, 1.10) -- (6.10, 1.10);
  \node[anchor=north] at (3.55, 0.88) {\footnotesize Timeout budget consumed};
  \draw[decorate, decoration={brace, mirror, amplitude=4pt}]
    (6.10, 1.10) -- (8.30, 1.10);
  \node[anchor=north] at (7.20, 0.88) {\footnotesize remaining budget};

  \draw[->] (-1.3, 0.42) -- (9.7, 0.42) node[anchor=south east] {time};

\end{tikzpicture}
}
  \caption{Timeout budget consumed under parallel capture.}
  \label{fig:overlap}
\end{figure}

\myparagraph{Existing thermal guard}
Samsung uses a static thermal guard to prevent Capture Timeout violations in deployed single-frame draft stages.
The guard uses Samsung's overheat level, an ordinal thermal indicator on a 0--6 scale, where larger values denote stronger throttling.
During pre-release validation, throttling at level~4 or higher repeatedly delayed draft sequences and caused Capture Timeout violations despite optimization.
The guard therefore skips optional single-frame stages at these levels.
However, the overheat level does not indicate a capture's waiting time or remaining draft work.
The guard therefore cannot directly assess whether the remaining timeout budget is sufficient for that work.

\myparagraph{Motivating experiment}
To examine the timeout risk of enabling lightweight multi-frame composition in the draft sequence, we compared two Portrait configurations on the Galaxy S26 Ultra.
We bypassed the thermal guard so that configured optional stages could execute regardless of the overheat level.
The optional stages are the multi-frame composition stage (\(M\)) and the deployed single-frame stages (collectively \(S\)).
Both \(M\) alone and \(M+S\) retain the same mandatory encoding and saving stages.
We tested 12MP and 24MP capture under normal memory conditions and under memory pressure.
For each condition in Table~\ref{tab:timeout_index}, we ran ten trials of up to 30 consecutive captures.
We chose 30 captures to extend testing beyond the earlier level-4 validation failure near the 20th consecutive capture.
The table reports the earliest timeout capture index across the ten trials, grouped by starting overheat level.

% Notes: docs/exhibits.md#tab_timeout_index
\providecolor{tmoS1}{RGB}{252,227,220} %
\providecolor{tmoS2}{RGB}{239,160,140} %
\providecolor{tmoS3}{RGB}{201,80,63}  %
\providecolor{tmoS4}{RGB}{143,36,32}  %

\newcolumntype{H}{>{\centering\arraybackslash}m{1.20cm}}
\newcolumntype{P}{>{\centering\arraybackslash}m{0.67cm}}

\providecommand{\tmoFirst}[2]{\cellcolor{#1}\makebox[\linewidth][r]{#2\hspace{1pt}}}
\providecommand{\tmoFirstW}[2]{\cellcolor{#1}\makebox[\linewidth][r]{\textcolor{white}{\textbf{#2}}\hspace{1pt}}}
\providecommand{\tmoNone}{\makebox[\linewidth][r]{--\hspace{1pt}}}

\begin{table}[t]
 \centering
 \caption{Capture index of the earliest timeout.}
 \label{tab:timeout_index}
 {\scriptsize\setlength{\tabcolsep}{1.3pt}%
 \begin{tabular}{@{}HPP|PP||PP|PP@{}}
 \toprule
 \multirow[c]{3}{=}[-10pt]{\makecell[c]{\textbf{Starting}\\\textbf{overheat}\\\textbf{level}}} &
 \multicolumn{4}{c||}{\raisebox{-3pt}{\textbf{Normal capture}}} &
 \multicolumn{4}{c@{}}{\raisebox{-3pt}{\textbf{Memory-pressure capture}}} \\[6pt] \cmidrule(lr){2-5}\cmidrule(l){6-9}
 &
 \multicolumn{2}{c|}{\raisebox{-3pt}{\textbf{12MP}}} &
 \multicolumn{2}{c||}{\raisebox{-3pt}{\textbf{24MP}}} &
 \multicolumn{2}{c|}{\raisebox{-3pt}{\textbf{12MP}}} &
 \multicolumn{2}{c@{}}{\raisebox{-3pt}{\textbf{24MP}}} \\[6pt] \cmidrule(lr){2-3}\cmidrule(lr){4-5}\cmidrule(lr){6-7}\cmidrule(l){8-9}
 & \raisebox{-2pt}{\textbf{M{+}S}} & \raisebox{-2pt}{\textbf{M}} &
 \raisebox{-2pt}{\textbf{M{+}S}} & \raisebox{-2pt}{\textbf{M}} &
 \raisebox{-2pt}{\textbf{M{+}S}} & \raisebox{-2pt}{\textbf{M}} &
 \raisebox{-2pt}{\textbf{M{+}S}} & \raisebox{-2pt}{\textbf{M}} \\[4pt]
 \midrule
 Lv0 &
 \tmoNone & \tmoNone & \tmoFirst{tmoS1}{29} & \tmoNone &
 \tmoNone & \tmoNone & \tmoFirst{tmoS1}{26} & \tmoNone \\
 Lv1 &
 \tmoNone & \tmoNone & \tmoFirst{tmoS1}{29} & \tmoNone &
 \tmoFirst{tmoS1}{24} & \tmoNone & \tmoFirst{tmoS2}{19} & \tmoNone \\
 Lv2 &
 \tmoFirst{tmoS1}{24} & \tmoNone & \tmoFirst{tmoS2}{18} & \tmoNone &
 \tmoFirst{tmoS2}{15} & \tmoFirst{tmoS1}{29} &
 \tmoFirst{tmoS2}{14} & \tmoFirst{tmoS1}{25} \\
 Lv3 &
 \tmoFirst{tmoS2}{13} & \tmoFirst{tmoS1}{21} &
 \tmoFirst{tmoS2}{12} & \tmoFirst{tmoS1}{20} &
 \tmoFirstW{tmoS3}{8} & \tmoFirst{tmoS2}{19} &
 \tmoFirstW{tmoS4}{5} & \tmoFirst{tmoS2}{13} \\
 \midrule
 Lv4 &
 \tmoFirstW{tmoS3}{8} & \tmoFirst{tmoS2}{16} &
 \tmoFirstW{tmoS3}{7} & \tmoFirst{tmoS2}{13} &
 \tmoFirstW{tmoS3}{6} & \tmoFirst{tmoS2}{13} &
 \tmoFirstW{tmoS4}{3} & \tmoFirstW{tmoS3}{7} \\
 Lv5 &
 \tmoFirstW{tmoS3}{8} & \tmoFirst{tmoS2}{12} &
 \tmoFirstW{tmoS3}{8} & \tmoFirst{tmoS2}{12} &
 \tmoFirstW{tmoS4}{4} & \tmoFirstW{tmoS3}{9} &
 \tmoFirstW{tmoS4}{4} & \tmoFirstW{tmoS3}{9} \\
 Lv6 &
 \tmoFirstW{tmoS3}{7} & \tmoFirst{tmoS2}{11} &
 \tmoFirstW{tmoS3}{7} & \tmoFirst{tmoS2}{11} &
 \tmoFirstW{tmoS4}{3} & \tmoFirstW{tmoS3}{6} &
 \tmoFirstW{tmoS4}{3} & \tmoFirstW{tmoS3}{6} \\
 \bottomrule
 \end{tabular}%
 }
 \par\vspace{3pt}
 {\scriptsize\begin{minipage}{0.933\columnwidth}
 \raggedright\setlength{\fboxsep}{1.8pt}%
 Timeout capture index:~%
 \colorbox{tmoS1}{\strut\,${\geq}\,20$\,}\colorbox{tmoS2}{\strut\,10--19\,}\colorbox{tmoS3}{\strut\,\textcolor{white}{6--9}\,}\colorbox{tmoS4}{\strut\,\textcolor{white}{${\leq}\,5$}\,}\par
 --: no timeout observed within 30 captures.\par
 \end{minipage}}
\end{table}


With the thermal guard bypassed, \(M+S\) timed out in 13 of the 16 conditions with starting levels below the guard's cutoff.
Every \(M+S\) trial starting at level~4 or higher nevertheless completed its first capture without a timeout, even though the guard would have skipped the optional stages for that capture.
Even at the same starting level, the earliest timeout varied with resolution, memory pressure, and draft configuration.
In conditions where \(M+S\) timed out, reducing optional work to \(M\) alone delayed the earliest timeout or resulted in no observed timeout within 30 captures.
Together, these observations motivate runtime control of optional-stage execution based on each capture's remaining timeout budget and draft workload, rather than the overheat level alone.